\section{Independent verification of the full Gram resonance
calculation}\label{independent-verification-of-the-full-gram-resonance-calculation}

25 September 2026. Independent mathematical verification of FGR1--FGR6
in
\href{https://github.com/KokunoYumeto/zeta-function-research-reader/blob/main/workbenches/splitzero-tandem/continuations/20260925-full-source-tensor-and-original-divisor/counterfactual/FULL_GRAM_RESONANCE_RETURN.md}{FULL\_GRAM\_RESONANCE\_RETURN.md}.
Proofs FGC1--FGC7 below also give a compact-uniform error estimate and
the exact sharp-window source remainder. These calculations take place
in the complex receiving spaces. Primitive \(Z_1/\tau\) acquires no
addition, parity, multiplicity or numerical weight.

\textbf{Verdict.} FGR1--FGR6 have the correct shift signs, Mellin
multiplier, Poisson density, Gram entries and multiplicity factors.
Their finite-set conclusion holds for every finite set of distinct
actual nontrivial zeta zeros, with all their multiplicities. No
critical-line hypothesis or reflection symmetry of that finite set is
required. Same-side endpoint vectors do not become orthogonal; use of
the complete Gram inverse is essential. The finite calculation does not
itself compare the projections with the full adelic cutoff or justify an
interchange with an exhaustion of all zeros.

\textbf{Human source and reading coverage.} Alain Connes,
\href{https://arxiv.org/abs/math/9811068v1}{\emph{Trace formula in
noncommutative geometry and the zeros of the Riemann zeta function},
arXiv:math/9811068v1}, §VIII, Lemma 3, equations (28)--(29). For this
verification the original author TeX was read at lines 2680--2720 of
\href{https://arxiv.org/abs/math/9811068v1}{connes98\_intake/math\_9811068v1\_author.tex}.
Lines 2696--2702 contain the individual asymptotic-orthogonality
assertion. The continuous-window calculation and proofs below are
programme derivations, not quotations or claims to have read the entire
source during this check.

\subsection{FGC1. The exact original Mellin receiver and its
sharp-window
remainder}\label{fgc1.-the-exact-original-mellin-receiver-and-its-sharp-window-remainder}

Let \(A\) be the programme space with seminorms

\[
 p_{N,j}(b)=\sup_{u>0}(u^N+u^{-N})|(u\partial_u)^j b(u)|,
 \qquad N,j\in\mathbb Z_{\geq0}.
\]

Define

\[
 (Ub)(x)=e^{x/2}b(e^x),\qquad
 (U^{-1}a)(u)=u^{-1/2}a(\log u).
\tag{FGC1.1}
\]

These are mutually inverse continuous linear maps between \(A\) and the
space \(\mathcal E\) of smooth functions for which
\(\sup_x e^{L|x|}|a^{(j)}(x)|<\infty\) for every real \(L\geq0\) and
every \(j\geq0\). To prove the assertion with all factors retained,
differentiate explicitly:

\[
 \partial_x^j(Ub)(x)
 =e^{x/2}\sum_{r=0}^j\binom jr(1/2)^{j-r}
 ((u\partial_u)^r b)(e^x),
\]

\[
 (u\partial_u)^j(U^{-1}a)(u)
 =u^{-1/2}\sum_{r=0}^j\binom jr(-1/2)^{j-r}a^{(r)}(\log u).
\tag{FGC1.2}
\]

Each requested target seminorm is bounded by a finite sum of source
seminorms, choosing an integer source exponent larger than the target
exponent plus \(1/2\). This proves both continuity statements and the
exact domains.

The original Mellin formula, without replacing the original zeta
function, is

\[
 M_0b(s)=\int_{\mathbb R}(Ub)(x)e^{(s-1/2)x}\,dx.
\tag{FGC1.3}
\]

Differentiating any number of times under this integral is valid locally
uniformly in \(s\), by the defining exponential decay of \(\mathcal E\).

For \(T>0\), the sharp window \(1_{[-T,T]}Ub\) corresponds under
(FGC1.1) to \(b_T(u)=1_{[e^{-T},e^T]}b(u)\). Generally \(b_T\notin A\),
because it can jump at the two endpoints. Its measurable compact-support
Mellin transform is nevertheless entire: on any compact set of \(s\),
all its differentiated integrands are integrable over the fixed finite
logarithmic interval. Thus the sharp-window formula has an exact
receiving domain; it is not implicitly an endomorphism of \(A\).

Let \(\rho\) be an actual zero and \(0\leq j<m_\rho\). The exact
identity is

\[
 \ell_{\rho,j,T}(Ub)
 =(M_0b)^{(j)}(\rho)
 -\int_{|x|>T}(Ub)(x)x^j e^{(\rho-1/2)x}\,dx.
\tag{FGC1.4}
\]

If \(b\in J=\Sigma S\), its full zero jets vanish, so

\[
 \ell_{\rho,j,T}(Ub)
 =-\int_{|x|>T}(Ub)(x)x^j e^{(\rho-1/2)x}\,dx.
\tag{FGC1.5}
\]

For every \(L>0\), the absolute value of the right side is
\(O_{b,\rho,j,L}(e^{-LT})\). Indeed use
\(|Ub(x)|\leq C e^{-(L+|\Re\rho-1/2|+2)|x|}\), absorb the fixed
polynomial into one exponential factor, and integrate the remaining
exponential tail. The remainder is small, but not asserted to be
identically zero. An identification of the finite constraint space with
a sharply truncated original source must account for this exact term.

\subsection{FGC2. Independence and the full finite
projection}\label{fgc2.-independence-and-the-full-finite-projection}

Write \(z_\rho=\rho-1/2\) as a coordinate difference, and let \(\Omega\)
be any finite set of distinct actual zeros with multiplicities
\(m_\rho\). Set \(M=\sum m_\rho\). The case \(M=0\) has zero projection
and every following identity is immediate, so suppose \(M>0\).

On \(L^2(\mathbb R,dx)\), with inner product conjugate linear in its
first argument, set

\[
 v_{\rho,j,T}(x)=1_{[-T,T]}(x)x^j e^{\overline{z_\rho}x}.
\]

Then \(\langle v_{\rho,j,T},\xi\rangle=\ell_{\rho,j,T}(\xi)\), which
verifies the conjugation, not only the span. To check independence,
assume \(\sum_\rho P_\rho(x)e^{\overline{z_\rho}x}=0\) on an interval,
with \(\deg P_\rho<m_\rho\). The entire exponential polynomial vanishes
identically. Fix \(\rho\) and apply
\(\prod_{\eta\ne\rho}(\partial_x-\overline{z_\eta})^{m_\eta}\). The
other summands vanish. On the surviving polynomial, each operator is
\(\partial_x+\overline{z_\rho}-\overline{z_\eta}\), invertible on
polynomials of degree below \(m_\rho\), since its matrix in the degree
basis is triangular with nonzero constant diagonal. Therefore
\(P_\rho=0\). This proves independence for all multiplicities.

Let \(V_T\) have these columns. Its Gram matrix is exactly

\[
 (G_T)_{(\rho,i),(\eta,j)}
 =\int_{-T}^T x^{i+j}e^{(z_\rho+\overline{z_\eta})x}\,dx.
\tag{FGC2.1}
\]

Independence makes \(G_T\) positive definite. Multiplication verifies
that \(P_T=V_TG_T^{-1}V_T^*\) is self-adjoint, idempotent, and identity
on the column span. For \(V(t)\xi(x)=\xi(x-t)\), finite-rank cyclicity
gives

\[
 \operatorname{Tr}(P_TV(t))
 =\operatorname{tr}(G_T^{-1}E_T(t)),
 \qquad E_T(t)=V_T^*V(t)V_T.
\tag{FGC2.2}
\]

A change of column basis by an invertible matrix \(C_T\) changes
\(G_T,E_T\) to \(C_T^*G_TC_T,C_T^*E_TC_T\). Their matrix trace in
(FGC2.2) is unchanged, since the new product is
\(C_T^{-1}G_T^{-1}E_TC_T\). Thus the endpoint and polynomial basis
changes below preserve the exact original projection, including every
cross term.

\subsection{FGC3. Endpoint and critical blocks, with cross-block
estimates}\label{fgc3.-endpoint-and-critical-blocks-with-cross-block-estimates}

For \(\Re z_\rho>0\), use the right-end basis

\[
 w^+_{\rho,j,T}(x)=1_{[-T,T]}(x)(T-x)^j
 e^{-\overline{z_\rho}(T-x)}.
\]

Its transformation from the original columns has nonzero degree diagonal
\((-1)^j e^{-\overline{z_\rho}T}\). In the coordinate \(y=T-x\), its
limit is \(y^je^{-\overline{z_\rho}y}\) on \(y\geq0\). For
\(\Re z_\rho<0\), use

\[
 w^-_{\rho,j,T}(x)=1_{[-T,T]}(x)(T+x)^j
 e^{\overline{z_\rho}(T+x)}.
\]

Its degree diagonal is \(e^{\overline{z_\rho}T}\), and its limit in
\(y=T+x\) is \(y^je^{\overline{z_\rho}y}\). Direct integration gives

\[
 G^+_{(\rho,i),(\eta,j)}
 =\frac{(i+j)!}{(z_\rho+\overline{z_\eta})^{i+j+1}},\qquad
 G^-_{(\rho,i),(\eta,j)}
 =\frac{(i+j)!}{(-z_\rho-\overline{z_\eta})^{i+j+1}}.
\tag{FGC3.1}
\]

The integral identity \(\int_0^\infty y^n e^{-cy}dy=n!c^{-n-1}\) for
\(\Re c>0\) follows first for \(n=0\) by its primitive and then
inductively by integration by parts. This proves every factor and sign
in (FGC3.1). The same exponential-polynomial independence argument
proves positivity of both limiting Gram matrices.

For \(\rho=1/2+i\gamma\), choose real polynomials \(p_j\) orthonormal
for \(dr/2\) on \([-1,1]\), and use

\[
 w^0_{\gamma,j,T}(x)=(2T)^{-1/2}1_{[-T,T]}(x)
 p_j(x/T)e^{-i\gamma x}.
\tag{FGC3.2}
\]

Its same-frequency Gram matrix is exactly identity: substitution
\(x=Tr\) gives \(\frac12\int p_i p_j\,dr=\delta_{ij}\). Distinct
critical frequencies give \(O(T^{-1})\), by one integration by parts in
the oscillatory integral; their frequency difference is nonzero because
the zeros in \(\Omega\) are distinct.

All cross-block estimates are uniform for \(|t|\leq L\) after inserting
one shift \(V(t)\), for any fixed \(L\). Here are their bounds with the
prerequisites explicit.

\begin{itemize}
\tightlist
\item
  An endpoint vector has absolute value bounded by a fixed polynomial in
  its endpoint distance times \(e^{-d y}\), where \(d=|\Re z_\rho|>0\).
  A critical vector is bounded by \(C T^{-1/2}\) on its window. Their
  overlap integral is therefore \(O(T^{-1/2})\); a shift changes
  endpoint distances by at most \(L\), which changes only the constant.
\item
  For opposite endpoints, on the overlapping windows their two endpoint
  distances have sum \(2T+t\) or \(2T-t\). Consequently the product is
  bounded by \(C_L(1+T)^K e^{-2d_*T}\), where \(d_*\) is the smaller of
  their two positive decay parameters. Integration adds a factor at most
  \(2T\). This tends to zero exponentially times a fixed polynomial.
\item
  For a same-side endpoint block, omitting its tail past \(2T\), with a
  bounded shift, produces another exponentially small polynomial tail.
\item
  For two critical vectors, the shifted shared interval becomes an
  interval in \([-1,1]\) with endpoints displaced by at most \(L/T\).
  Their coefficient integrand is a bounded polynomial in \(r,r-t/T\),
  multiplied by \(e^{i(\gamma-\gamma')Tr}\) and the exact factor
  \(e^{i\gamma't}\). With equal frequencies, its difference from the
  unshifted coefficient is \(O_L(T^{-1})\), both on the shared interval
  and on the lost intervals. With unequal frequencies, integration by
  parts gives \(O_L(T^{-1})\).
\end{itemize}

It follows that the Gram matrix in this basis tends to

\[
 B=G^+\oplus G^-\oplus\bigoplus_{\gamma}I_{m_{1/2+i\gamma}}.
\tag{FGC3.3}
\]

Empty blocks are omitted. Its inverse exists, and continuity of
inversion in finite matrices follows, for example, from the convergent
Neumann series for \((I+B^{-1}(G_T-B))^{-1}\) once the operator norm of
its argument is below one. There is no orthogonality assumption inside
either endpoint block.

For two simple right-end vectors their individual unit-vector inner
product tends to

\[
 \frac{2\sqrt{(\beta_\rho-1/2)(\beta_\eta-1/2)}}
 {z_\rho+\overline{z_\eta}},
\tag{FGC3.4}
\]

which is nonzero. In the discrete finite-window example with columns
\(2^n,3^n\), direct geometric sums give the limit \(\sqrt{24}/5\). Thus
the correction to individual orthogonality in FGR2 is independently
confirmed.

\subsection{FGC4. The trace, all multiplicities, and a quantitative
extension}\label{fgc4.-the-trace-all-multiplicities-and-a-quantitative-extension}

At the right endpoint \(y=T-x\), the shift \(V(t)\) becomes
\(f(y)\mapsto f(y+t)\). For \(t\geq0\), the limiting half-line
compression preserves the finite endpoint span and acts by

\[
 y^j e^{-\overline{z_\rho}y}
 \longmapsto e^{-\overline{z_\rho}t}
 \sum_{r=0}^j\binom jr t^{j-r}y^r e^{-\overline{z_\rho}y}.
\tag{FGC4.1}
\]

Let \(C_+(t)\) be this complete triangular coefficient matrix. The
limiting matrix of inner products is \(G^+C_+(t)\), so the full Gram
trace is \(\operatorname{tr}C_+(t)\), exactly
\(\sum_{\Re z_\rho>0}m_\rho e^{-\overline{z_\rho}t}\). The off-diagonal
nilpotent entries remain present in (FGC4.1); their diagonal sum is
zero.

For negative \(t\), the compression is the adjoint of the positive-time
compression. Indeed extension by zero followed by compression commutes
with taking the adjoint, and \(V(-t)=V(t)^*\). Its trace is the complex
conjugate. The right-end contribution at every real \(t\) is therefore
\(\sum_{\Re z_\rho>0}m_\rho e^{-(\beta-1/2)|t|}e^{i\gamma t}\).

At the left endpoint \(y=T+x\), the shift becomes
\(f(y)\mapsto f(y-t)\). For \(t\leq0\), its restriction is triangular
with diagonal \(e^{-\overline{z_\rho}t}\); for positive time use its
adjoint. This yields
\(\sum_{\Re z_\rho<0}m_\rho e^{-(1/2-\beta)|t|}e^{i\gamma t}\).

For a critical column the limiting shift matrix is \(e^{i\gamma t}I\).
This is the positive sign: the column oscillation is \(e^{-i\gamma x}\),
so replacing \(x\) by \(x-t\) multiplies it by \(e^{i\gamma t}\).
Combining the blocks proves

\[
 \lim_{T\to\infty}\operatorname{Tr}(P_TV(t))
 =\sum_{\rho\in\Omega}m_\rho
 e^{-|\beta-1/2||t|}e^{i\gamma t}.
\tag{FGC4.2}
\]

There is also the stronger bound

\[
 \sup_{|t|\leq L}\left|
 \operatorname{Tr}(P_TV(t))-
 \sum_{\rho\in\Omega}m_\rho e^{-|\beta-1/2||t|}e^{i\gamma t}
 \right|=O_{\Omega,L}(T^{-1}).
\tag{FGC4.3}
\]

To prove it without ignoring the mixed entries, write \(G_T=B+R_T\) and
\(E_T(t)=E_\infty(t)+S_T(t)\) in the basis of FGC3. Here \(E_\infty\) is
block diagonal, its endpoint blocks are the complete compressed-shift
matrices of inner products, and its critical blocks are
\(e^{i\gamma t}I\). The cross-block bounds prove
\(\|R_T\|,\sup_{|t|\leq L}\|S_T(t)\|=O(T^{-1/2})\). The diagonal blocks
of each error are \(O(T^{-1})\), or exponentially smaller for endpoints.
All constants are finite because \(\Omega\) is fixed.

Expand the inverse with its retained quadratic remainder:

\[
 (B+R_T)^{-1}=B^{-1}-B^{-1}R_TB^{-1}+O(T^{-1}).
\]

The trace difference is

\[
 \operatorname{tr}(B^{-1}S_T(t))
 -\operatorname{tr}(B^{-1}R_TB^{-1}E_\infty(t))
 +O(T^{-1}).
\tag{FGC4.4}
\]

In both displayed traces only the diagonal blocks of the errors
contribute, because the other factors are block diagonal. They are
\(O(T^{-1})\). The omitted product \(R_TS_T\) and the inverse remainder
are also \(O(T^{-1})\). This proves (FGC4.3). In particular the mixed
\(T^{-1/2}\) entries were included, and their products explain why the
trace has a stronger rate than the column convergence.

The constants are not claimed uniform as \(\Omega\) grows, as an
off-critical real part approaches \(1/2\), or as distinct frequencies
approach each other. Thus (FGC4.3) does not authorize an
infinite-zero/window interchange.

\subsection{FGC5. Integrated tests and the exact Poisson
factor}\label{fgc5.-integrated-tests-and-the-exact-poisson-factor}

For \(b\in A\), put \(a_b(t)=e^{t/2}b(e^t)\). By FGC1 it belongs to
\(L^1(\mathbb R)\). The strong operator integral
\(V(a_b)=\int a_b(t)V(t)dt\) has norm at most \(\|a_b\|_1\). Rank \(M\)
and unitarity give \(|\operatorname{Tr}(P_TV(t))|\leq M\). Therefore
dominated convergence, without any compact-support assumption on
\(a_b\), yields

\[
 \lim_{T\to\infty}\operatorname{Tr}(P_TV(a_b))
 =\sum_{\rho\in\Omega}m_\rho\int_{\mathbb R}
 a_b(t)e^{-|\beta-1/2||t|}e^{i\gamma t}\,dt.
\tag{FGC5.1}
\]

For \(d>0\), residue integration of
\(d\,e^{iyt}/[\pi(d^2+(y-\gamma)^2)]\) gives

\[
 \int_{\mathbb R}\frac{d\,e^{iyt}}{\pi(d^2+(y-\gamma)^2)}dy
 =e^{-d|t|}e^{i\gamma t}.
\tag{FGC5.2}
\]

For positive time the upper pole \(\gamma+id\) contributes
\(e^{i\gamma t-dt}\); for negative time the lower pole and clockwise
orientation give \(e^{i\gamma t+dt}\). At zero the integral of the
density is one. Fubini is justified by this mass and \(\|a_b\|_1\), so
an off-critical term is exactly

\[
 m_\rho\int_{\mathbb R}M_0b(1/2+iy)
 \frac{|\beta-1/2|}{\pi((\beta-1/2)^2+(y-\gamma)^2)}dy.
\tag{FGC5.3}
\]

There is no additional factor of \(2\pi\). The critical term is
\(m_\rho M_0b(\rho)\). Both follow from the same original Mellin
convention (FGC1.3).

In fact (FGC5.1) holds for every \(a\in L^1\) in place of \(a_b\), with
the same finite sum of kernels. If \(a^*(t)=\overline{a(-t)}\), Fubini
gives \(V(a*a^*)=V(a)V(a)^*\). Its finite-rank compressed trace is a sum
of squared norms, hence nonnegative. Equivalently the limiting positive
measure on the frequency axis is

\[
 \sum_{\substack{\rho\in\Omega\\\beta=1/2}}m_\rho\delta_\gamma
 +\sum_{\substack{\rho\in\Omega\\\beta\ne1/2}}
 m_\rho\frac{|\beta-1/2|\,dy}
 {\pi((\beta-1/2)^2+(y-\gamma)^2)}.
\tag{FGC5.4}
\]

Its total mass is precisely \(M\). This positivity statement is about
the calculated finite receiving projection.

\subsection{FGC6. The original arithmetic discrepancy and reflected
pair}\label{fgc6.-the-original-arithmetic-discrepancy-and-reflected-pair}

In the original Mellin formula the contribution of \(\rho\) is

\[
 m_\rho M_0b(\rho)=m_\rho\int a_b(t)
 e^{(\beta-1/2)t}e^{i\gamma t}dt.
\]

Subtracting (FGC5.1) gives exactly

\[
 \sum_{\rho\in\Omega}m_\rho\int a_b(t)e^{i\gamma t}
 \left(e^{(\beta-1/2)t}-e^{-|\beta-1/2||t|}\right)dt.
\tag{FGC6.1}
\]

If \(\Omega\) contains both \(\beta+i\gamma\) and \(1-\beta+i\gamma\),
each of multiplicity \(m\), and \(d=|\beta-1/2|>0\), the sum of their
two exponential differences is

\[
 e^{dt}+e^{-dt}-2e^{-d|t|}=2\sinh(d|t|).
\]

The pair contributes \(2m\int a_b(t)e^{i\gamma t}\sinh(d|t|)dt\). This
verifies the sign and factor in FGR6.3. Reflection symmetry is used only
in this pair identity, not in the projection or its convergence theorem.

\subsection{FGC7. What this check proves and what the next comparison
must
retain}\label{fgc7.-what-this-check-proves-and-what-the-next-comparison-must-retain}

No error was found in FGR1--FGR6. The independent proof establishes the
finite trace limit with all nilpotent orders and arbitrary finite
configurations of zeros, the exact original Mellin return, the
positivity measure, and the original arithmetic discrepancy. The
quantitative extension (FGC4.3) and the exact source-window remainder
(FGC1.5) strengthen the finite result.

For the next actual source comparison, two concrete terms are now
explicit: the sharp-window Mellin-jet tail (FGC1.5), and the
arithmetic-minus-harmonic expression (FGC6.1). The former tends to zero
faster than any fixed exponential for a fixed source test and finite
zero set. The latter is a different trace kernel and is not thereby
zero. Testing that second expression on actual elements of \(J\), with
the full original arithmetic formula, is therefore a mathematically
specified next calculation. This check neither assumes an off-critical
zero nor proves that no such zero exists.

\subsection{Publication source and dependency
links}\label{publication-source-and-dependency-links}

The source geometry is Alain Connes and Caterina Consani, \emph{Schemes
over F1 and zeta functions},
\href{https://arxiv.org/abs/0903.2024v3}{arXiv:0903.2024v3, §5}. The
weight-control comparison is Pierre Deligne, \emph{La conjecture de
Weil. II}, Publications mathematiques de l'IHES52 (1980),137-252,
\href{https://www.numdam.org/item/PMIHES_1980__52__137_0/}{§§3.3.11
and3.6}. Deligne material was read through the identified French
transcription and separately recorded peer source-page checks, not
Deligne-authored TeX. These sources are not asserted to contain the new
programme derivations. The
\href{https://github.com/KokunoYumeto/zeta-function-research-reader/blob/36a82ecca98addb8a98b48204e349737856c7223/workbenches/splitzero-tandem/continuations/20260924-cohomology-weight-and-character-lifts/BUILD_AND_REVIEW.md}{preceding
complete proof and reading record} records inherited source versions and
actual inspection limits. The investigator's corrected construction
remains attributed in the original text above.
